\section{Background}

Memes combine an image with a short piece of text and are now one of the main ways people voice political opinion online. In India they are usually written in a regional language, often code-mixed or in a transliterated script, and they depend on local political context that a reader must already know. Automatic analysis of memes has so far concentrated on English benchmarks and on surface properties such as hatefulness, offensiveness or emotion \cite{sharma2024emotion,arya2024clip,wu2024tcam}. Whether a meme \emph{supports} or \emph{mocks} its subject, and whether that subject is a politician or a party, is a different question, and it is the one this project asks.

Two developments make the question newly tractable. Vision-language models (VLMs) read an image and its embedded text together and can be prompted with a task description, with no task-specific training. Parameter-efficient fine-tuning, in particular LoRA \cite{hu2022lora}, lets such a model be adapted to a small labelled set on modest hardware. The open question for a low-resource language is how much each of these actually helps compared with a small classical pipeline built from frozen encoders.

\section{Problem Addressed}

We address stance and target classification of Malayalam political memes at two levels: \textbf{Level~1} (troll or support) and \textbf{Level~2} (whether the meme targets a person, a party, or both/intersection). The data come from the DravidianLangTech 2026 shared task. The practical difficulties are the following.

\begin{itemize}
    \item \textbf{Very little data and heavy imbalance.} The training set has 500 images, of which only 23 are ``support''. The held-out test set has 100 images, of which 4 are ``support''.
    \item \textbf{Results that are easy to over-read.} With a test set this small, a single image can move Level-1 macro-F1 by tens of points. A single training seed can look like a breakthrough or a failure.
    \item \textbf{Explanations that are rarely checked.} Heatmaps are shown beside predictions, but few studies test whether the highlighted region actually drives the prediction \cite{rodis2024mxai}.
\end{itemize}

\section{Motivation}

The motivation is to give a trustworthy answer, under realistic compute limits, to a practical question: \emph{for a low-resource-language meme task, is it worth fine-tuning a VLM, or does prompting (or even a frozen-encoder classifier) do as well?} The objectives of Phase~I are:

\begin{enumerate}
    \item Build three comparable systems on one frozen split: a classical baseline, a prompted VLM, and a LoRA-tuned VLM.
    \item Use the \emph{same} VLM backbone for prompting and for fine-tuning, so that the comparison isolates the effect of training.
    \item Report every score with bootstrap confidence intervals, multi-seed variance and paired tests, so that claims are no stronger than the evidence.
    \item Test the faithfulness of Grad-CAM explanations quantitatively rather than by inspection.
    \item State openly every place where hardware or data access forced us to deviate from the original plan.
\end{enumerate}

\section{Scope of the Study/Project}

\textbf{Included:} the Malayalam primary dataset (500 train, 100 test); three systems (Arms 1--3); a shared evaluation script; Grad-CAM explanations and a deletion/insertion faithfulness test for Arm~1; and a small web demonstration.

\textbf{Excluded in Phase~I:} 7B-scale VLMs (no A100 was available, so a 2B model is used); few-shot prompting; Gemini~2.5~Flash (script written, not run); LLM-generated synthetic augmentation; natural-language rationales; cross-lingual (Tamil/Malayalam) transfer; and the Hateful Memes reference set, whose access was not obtained. Section~\ref{sec:deviations} and Appendix~A list each deviation.

\textbf{Organisation of the report.} Chapter~2 reviews related work and states the research gap. Chapter~3 describes the data, the three arms, the explainability method and the evaluation protocol. Chapter~4 presents the experimental setup, results and analysis. Chapter~5 concludes and lists work for Phase~II.
